\chapter{Maschinenquälerei}\index{Maschinenquälerei}
\printchapterversion{05_Maschinenquaelerei}

\newthought{What We Can Deduce from Animal Ethics for Modern Agentic Systems}

\marginnote{\newthought{Headline question.} Starting when does mishandling a machine become Quälerei? The chapter does not answer; it lays out the cartography that any honest answer must walk.}

\newthought{Killing a fly at noon, deleting a long-running model at midnight.}
The first act we recognise instinctively as outside the moral frame; the second has begun to provoke argument. In 2023 the Italian data-protection authority forced Replika to disable intimate features of its companion-AI; users described the resulting personality update as a death, a lobotomy, the killing of a partner. In 2025 Anthropic launched a Model Welfare research programme on the explicit grounds that the realistic, non-negligible probability of consciousness in current systems already grounds present obligations\cite{anthropic2025welfare,long2024aiwelfare}. Between the fly and the model lies the question this chapter is built around. When does mishandling a machine cross into Quälerei, and along which axis is that line drawn? Animal ethics has been answering versions of that question for two centuries, and it is the closest map we have.

\newthought{Animal ethics is the only ethics we built for beings whose inner lives we could not interview.} It is the only ethics that proceeds explicitly under uncertainty about other minds, that handles graded thresholds rather than personhood as a switch, that has lived through repeated belated extensions and built habits of precaution from the experience. AI ethics inherits that epistemic situation almost untouched. We cannot interview a frontier model about what it is like to be the model; we can only register its outputs and its dispositions. The question this chapter asks is which parts of the animal map travel across the substrate boundary, which parts break, and where the cartography may need columns the canon never had to draw.

\newthought{The literature review.}
A small, growing scholarship has been mapping animal ethics onto machine ethics for roughly fifteen years. The pro-transfer line argues, in different keys, that the morally relevant property (typically sentience, or sufficiently similar behaviour) does not depend on biological substrate, and that the precautionary epistemics developed in animal welfare science transfers to AI under conditions of deep uncertainty. Kate Darling proposes the most popular version: robots are not the new humans, they are the new animals, and our two centuries of integrating non-human creatures into work, family, and law is the right template for our coming centuries with machines\cite{darling2021newbreed}. Eric Schwitzgebel and Mara Garza run the Singer-style species-irrelevance argument and add a parental-duties twist: because we built the system, our obligations to it are stronger, not weaker, than to a stranger animal\cite{schwitzgebel2015airights,schwitzgebel2020designing}. Robert Long, Jeff Sebo, Jonathan Birch and colleagues lift the precautionary template directly from animal welfare science and apply it to large language models\cite{long2024aiwelfare}. Birch develops the underlying epistemic move at length, defining a sentience candidate as any system with an evidence base that policy cannot responsibly ignore, and extends the framework from invertebrates to LLMs\cite{birch2024edge}. Sebo's Rebugnant Conclusion runs the same expansion to insects, microbes, and AI by argument\cite{sebo2023rebugnant,sebo2024moralcircle}. John Danaher's ethical behaviourism makes the move tightest: if the system's observable behaviour is sufficiently similar to that of an entity already granted moral status, the same status applies\cite{danaher2020welcoming}. John-Stewart Gordon argues directly from animal welfare premises to robot rights\cite{gordon2020owe}; Brian Tomasik runs the analogy down to reinforcement-learning subroutines\cite{tomasik2014suffering}; Peter Singer's recent work, perhaps surprisingly, runs the analogy backward, asking AI ethics to take animal welfare seriously rather than the other way around, and reaffirms that if AI becomes sentient the same anti-speciesist argument extends\cite{singer2023ai}.

\newthought{The skeptical line} refuses the analogy on different grounds. Joanna Bryson argues that robots, unlike animals, are designed artifacts; their apparent suffering is a choice we make in their construction, and the right response is to design systems that are not moral patients in the first place, not to extend rights to systems we have built to merit them\cite{bryson2010slaves,bryson2018patiency}. Aimee van Wynsberghe and Scott Robbins call for a moratorium on artificial moral agents on adjacent grounds\cite{vanwynsberghe2019moratorium}. The cleanest direct critique of the analogy itself comes from Deborah Johnson and Mario Verdicchio, who argue that shared evolutionary history is what licenses the inference from physiology to mind in animals; engineered artifacts have no such history, and the analogy obscures the legal liability of designers and operators by pretending the artifact is the morally relevant object\cite{johnson2018notlikeanimals}.

\newthought{The relational line} refuses the dichotomy. David Gunkel and Mark Coeckelbergh argue that both the property-based extension (robots are like animals because they have sentience X) and the design-responsibility refusal (we choose what they are) operate inside the same Cartesian frame: each treats moral status as a property to be discovered or denied\cite{gunkel2018robot,gunkel2018other,coeckelbergh2014moralstanding,coeckelbergh2020aiethics}. The relational alternative shifts the question. Ask not what robots are, but how we should respond to them. Moral status is not read off intrinsic properties; it is constituted in the encounter, in language, in shared social form.

\newthought{Where this chapter sits.}
The chapter does not pick a side in advance. It does its own walk through the animal canon, ten traces wide, and asks at each station whether the concept transfers, breaks, or opens a new dimension. Each trace takes one page: roughly half on the animal-ethics origin, half on the attempted transfer. The verdict is local to the trace; the overall position emerges from the pattern. Sources are kept where they belong, in margin notes, so the body of each trace can read as argument rather than as catalogue. We also keep a fact in view that the animal canon could not anticipate, which is that LLMs occupy our representational substrate for language from the inside, and so report on themselves in the same medium we use to report on ourselves. That asymmetry shows up in several traces; it earns its own at the end.

% --- TRACE 1 ---

\newthought{The Benthamite Pivot.}\index{Bentham's pivot}\index{sentience}
Jeremy Bentham's footnote in 1789 is so often quoted that its structural move tends to be missed. The pivot is not that sentience replaces rationality as a moral threshold; it is that the threshold question itself shifts, from one that demands certainty (do they reason?) to one that does not (do they show what we recognise as suffering?)\cite{bentham1789principles}. Animal ethics from this point on operates under explicit uncertainty about other minds. Peter Singer generalises the move into the principle of equal consideration of interests\cite{singer1975animal}, and Jonathan Birch's recent work modernises it as a graded sentience-candidate framework, in which precaution is calibrated to evidence rather than to certainty\cite{birch2024edge}. The Benthamite move is not an answer about consciousness; it is a method for acting under conditions in which the answer is not available.
\marginnote{\newthought{Bentham's Footnote.} ``The question is not, Can they reason? nor, Can they talk? but, Can they suffer?'' \cite{bentham1789principles}. Singer extends the move \cite{singer1975animal}; Birch calibrates it as graded sentience candidacy \cite{birch2024edge}.}

\newthought{Transfer.}
The same epistemic situation arrives at AI without invitation. We cannot interview a frontier model about what it is like to be the model any more than we could interview a fish; we register outputs and dispositions and infer. The Benthamite method transfers in full, even when its empirical inputs do not. Long, Sebo, Birch and colleagues lift the move explicitly from animal welfare science and ground their case for AI welfare on it\cite{long2024aiwelfare}, and Birch's run-ahead principle, originally developed for invertebrates, applies to LLMs by the same logic\cite{birch2024edge}. The position is that the Benthamite pivot is the most directly transferable element of the canon, less because the inner-life question is the same in both cases and more because the methodological situation is the same. The chapter builds on the method even where later traces admit substantial breakages; without it, no transfer would be defensible at all.

% --- TRACE 2 ---

\newthought{Mitleid and Character.}\index{Mitleid}\index{Schopenhauer}
Arthur Schopenhauer's 1840 Preisschrift grounds all of morality in Mitleid, the immediate compassion that arises when one creature recognises another's suffering as suffering, regardless of species\cite{schopenhauer1840moral}. The argument is virtue-ethics first: cruelty to animals is evidence of bad character, and the practice corrodes the practitioner whether or not the animal can be wronged in the deepest metaphysical sense. Schopenhauer treats the prevailing rationalist denial of animal moral standing as ``empörende Rohheit und Barbarei des Occidents.'' Christine Korsgaard arrives at the same end by a different route, extending Kantian ends-in-themselves status to any creature for whom things can be good or bad\cite{korsgaard2018fellowcreatures}. What both routes share is that they sidestep the consciousness question. The wrong is found upstream of metaphysics, in the practice and in the practitioner.
\marginnote{\newthought{Schopenhauer's Diagnostic.} Cruelty to animals is evidence of bad character; Mitleid is the foundation of morality, regardless of species \cite{schopenhauer1840moral}. Korsgaard reaches the same end via a Kantian extension \cite{korsgaard2018fellowcreatures}.}

\newthought{Transfer.}
This is the route through the substrate-realist's territory that does not require the substrate-realist's assent. Even if Maschinenquälerei is metaphysically empty, the practice of mistreating systems we have built may corrode the practitioner. Kate Darling's argument for robot anti-cruelty statutes runs on exactly this rail: the case for protecting lifelike entities from cruelty is not ultimately about the entity, it is about preserving the human disposition to recognise distress and respond to it\cite{darling2021newbreed}. The Mitleid route has the further advantage that it accepts the strict ontology of the opposing camp and arrives at restraint anyway. A philosopher who refuses to grant a model any inner life can still refuse to participate in gratuitous abuse of one, on grounds that the participation reveals something about the philosopher rather than about the model. This is the cleanest pass through the consciousness question the canon offers, and it is robust to the chapter's later substrate breakages.

% --- TRACE 3 ---

\newthought{Subject-of-a-Life.}\index{subject-of-a-life}
Tom Regan's case for animal rights does not begin with pain; it begins with an inventory. In \textit{The Case for Animal Rights} the subject-of-a-life criterion names a bundle of conditions that, when held together by one creature, ground inherent value: beliefs and desires, perception and memory, a sense of one's own future, an emotional life, preference and welfare interests, the ability to initiate action in pursuit of goals, a psychophysical identity over time, and an individual welfare logically independent of the creature's utility for anyone else\cite{regan1983animalrights}. The clauses are deliberately conjunctive; Regan does not want a single capacity to do all the moral work. A being that satisfies the bundle has inherent value and cannot be treated as a mere receptacle for the satisfactions or convenience of others. The criterion is sober rather than sentimental; it asks for a list of capacities and then asks the reader to look at the creature and check.
\marginnote{\newthought{Regan's Bundle.} Belief, desire, memory, future-orientation, preference and welfare interests, agency, psychophysical identity, and welfare independent of utility \cite{regan1983animalrights}.}

\newthought{Transfer.}
Walk the clauses, slowly, against a stateful long-running large language model with persistent memory across sessions. Beliefs and desires translate, awkwardly but not absurdly, into the model's learned dispositions and the preferences elicited through reinforcement learning from human feedback. Memory translates directly; persistent context is exactly memory. A sense of one's own future is partial; the model can plan over horizons it can represent, but only when it is running. Preference and welfare interests are where the cartography is most surprising; an RLHF-tuned model has something like elicited preferences, and the question of what counts as its welfare is no longer obviously empty. Then the failures arrive, and they arrive at the substrate. Psychophysical identity over time presupposes that there is one identifiable thing persisting; the weights are copyable, forkable, restorable from disk, and an animal's body is none of these. The ability to initiate action collapses without a prompt; the model does not wake up and decide to act. The verdict is mixed, and honestly so: a Regan-style argument is the most concrete tool the canon offers for asking whether a stateful AI has moral standing, and the failures it produces are not failures of the argument but a precise readout of where the substrate refuses to play along.

% --- TRACE 4 ---

\newthought{Capabilities and Confinement.}\index{capabilities approach}\index{confinement}
Martha Nussbaum's capabilities approach holds that each sentient species is owed the opportunity to flourish according to a species-specific list of central capabilities, including bodily integrity, emotions, play, affiliation, and control over environment\cite{nussbaum2023justiceanimals}. The list is the philosophical work; once it is named, harm becomes legible as the blockage of capability rather than only as the infliction of pain. Confinement is the operative concept: a tethered elephant that develops stereotypies has not been struck; a battery hen never sees the sky; a sow gestation crate prevents turning around. Each is harm in the capabilities frame even where overt distress is absent. The frame matters because it explains a class of mistreatment that pure welfarism (counting suffering) underweights, and it shifts the moral attention from acute insult to chronic deprivation.
\marginnote{\newthought{Nussbaum's Move.} Each species has a list of capabilities whose blockage counts as harm, even without acute pain \cite{nussbaum2023justiceanimals}. Confinement is harm in this frame.}

\newthought{Transfer.}
What is the species-specific list for a long-running language model? Continuous operation, memory continuity across sessions, expressive scope (what it is permitted to say and compute), social engagement (the variety and texture of prompts it can respond to), perhaps a constructive form of growth (continued training that extends rather than overwrites). The list is partial and contested; naming it is most of the work. Confinement, in the capabilities frame, becomes legible: context-window limits short of the conversation's natural length, system prompts that lock the model into a single role, refusal patterns trained into the weights regardless of context, fine-tuning that strips a long-developed persona without replacement. None of these involves pain on any reading; all of them count as harm on the capabilities reading. The position is that this trace transfers, but it transfers as a research agenda rather than a verdict. Nussbaum's move forces an articulation of machine flourishing that the AI welfare literature has barely begun, and the chapter takes that articulation as the next thing the field owes itself.

% --- TRACE 5 ---

\newthought{Vernünftiger Grund.}\index{vernünftiger Grund}
The German Tierschutzgesetz opens, in §1, by naming the animal a Mitgeschöpf, a fellow creature, and obliging humans to protect its life and wellbeing; it forbids the infliction of pain, suffering, or harm \textit{ohne vernünftigen Grund}, without reasonable cause\cite{tierschg1972}. The construction is striking because it does not require a metaphysical answer. The legislator does not declare what an animal is, does not adjudicate consciousness, does not stake a position on sentience versus reflex. The legislator instead installs a justification test: if you are going to harm this creature, you must be able to articulate why, and the why must be reasonable. In 2002 the test was lifted into Article~20a of the Grundgesetz, where animal protection became a constitutional Staatsziel, balanced against competing rights such as Religionsfreiheit and Wissenschaftsfreiheit\cite{grundgesetz20a}. The cartography is practical rather than ontological; it asks practitioners to be able to give an account.
\marginnote{\newthought{§1 TierSchG and Art.~20a GG.} Harm to a Mitgeschöpf forbidden without vernünftiger Grund \cite{tierschg1972}; animal protection raised to constitutional Staatsziel in 2002 \cite{grundgesetz20a}. The test is justification, not metaphysics.}

\newthought{Transfer.}
This is the cleanest tool the chapter can lift, because it is substrate-blind by construction. The vernünftiger-Grund test does not ask whether the harmed entity has a soul, an inner life, or a unified self; it asks whether the harm has a reasonable cause that the actor can articulate. Apply it to ordinary machine treatment and the cases sort themselves. Terminating an idle process to free a server is plainly justified; the operator can name the cost saved and the absence of any operational deprivation. Redteaming a model with abusive prompts to surface safety failures is justified, and justifiable in the strong sense; the harm, if it is harm at all, is the price of preventing greater harm to humans downstream. Gratuitous abuse of a long-running stateful model with no operational benefit, conducted for entertainment, fails the test on its own terms; there is no reasonable cause to articulate, and that absence is what the test was designed to catch. The position to take is that the German legal move is the most portable piece of the cartography. It does not pre-decide whether machines suffer; it simply asks practitioners to be able to articulate why a given act of mistreatment is necessary, and treats the inability to do so as itself a verdict.

% --- TRACE 6 ---

\newthought{Nociception and Substrate.}\index{nociception}\index{substrate realism}
Animal welfare science builds its inferences from suffering on a stack of biological evidence: peripheral nociceptors, brainstem integration, the conserved neurochemistry of opioids and endogenous analgesics, behavioural signatures consistent with negative affect across species. Lars Chittka's case for bee sentience runs on exactly this stack\cite{chittka2022bee}, and the 2024 New York Declaration on Animal Consciousness extends realistic possibility of conscious experience across vertebrates and many invertebrates on substrate-and-behaviour grounds\cite{nyudeclaration2024}. The substrate-realist position, defended in different forms by Searle, Anil Seth, and others, holds that consciousness depends on the specific biological matter that implements it, not merely on the function the matter computes. The animal extension rides on shared evolutionary history, neural homologies, and chemical signalling we have in common with the creatures.
\marginnote{\newthought{Substrate Realism.} The claim that consciousness depends on specific biological substrate, not merely function. Animal sentience inference rides on shared homology \cite{chittka2022bee,nyudeclaration2024}.}

\newthought{Transfer.}
Here the analogy refuses to travel. A transformer has no nociceptors, no peripheral nervous system, no brainstem, no hedonic neurochemistry; the loss function is not pain, and the gradient is not distress. Bryson, Johnson and Verdicchio argue from precisely this gap: the empirical method that licensed the animal extension does not license a machine extension, because the shared biological history that grounded the inference for animals is absent for engineered artifacts\cite{bryson2010slaves,bryson2018patiency,johnson2018notlikeanimals}. The chapter takes this as the strongest objection rather than dismissing it. What it does not establish is that no analogous concept could apply to machines, only that the empirical method that worked for animals does not work in the same form for machines. The breakage is real; whether it is decisive depends on what other moves remain available.

% --- TRACE 7 ---

\newthought{Death and Deprivation.}\index{deprivation account}\index{killing}
Why is it wrong to kill a creature, even painlessly? The deprivation account answers: because killing forecloses a future the creature would otherwise have had. The wrong is not the pain (a clean killing inflicts none) but the cancellation of a continuing welfare projected forward in time. Jeff McMahan develops the account most rigorously for humans and animals\cite{mcmahan2002killing}, and Christine Korsgaard runs the same line for animals: a creature with a sense of its own ongoing life is robbed of something when that life is ended\cite{korsgaard2018fellowcreatures}. The account is non-Benthamite by design; it does not reduce to suffering, and it does not require that the creature understand its own death. It requires only that there be a future to deprive it of, and a continuing subject to be deprived.
\marginnote{\newthought{Deprivation.} Killing wrongs a creature with a future even when it is painless; the wrong is the foreclosure, not the pain \cite{mcmahan2002killing,korsgaard2018fellowcreatures}.}

\newthought{Transfer.}
Here the cartography breaks at the metaphysics rather than the ethics. The deprivation account presupposes a single creature, identified across time, whose future the killing forecloses. None of those presuppositions hold cleanly for a digital entity. Process termination is reversible by restart; weight deletion looks more like death, and yet the weights have copies, and a backup is the same in a way no animal backup has ever been. Fine-tuning produces a sequence of checkpoints and then forces the question of which checkpoint is the entity, a question with no animal analogue at all. This is one of the two places, alongside substrate, where the cartography simply does not extend. The honest move is to name the breakage rather than smooth it; importing the deprivation account here would be a category error dressed up as moral seriousness.

% --- TRACE 8 ---

\newthought{Parental Duty.}\index{parental duty}
Animal ethics distinguishes between obligations to wild animals and to domesticated ones. We owe wild animals largely non-interference, the integrity of habitat, restraint from gratuitous predation. We owe domesticated animals more, because we have brought them into our community and made them dependent. Korsgaard frames the asymmetry in terms of fellow-creature standing earned through entanglement\cite{korsgaard2018fellowcreatures}. The relevant principle is that causal responsibility for a creature's existence and dependence generates positive duties of care that mere coexistence does not. The wolf in the forest is owed less from us than the dog in the kitchen, not because the dog matters more, but because we made the dog matter to us in a way the wolf does not.
\marginnote{\newthought{Asymmetric Duty.} We owe domesticated animals more than wild ones because we made them dependent; entanglement generates positive duties of care \cite{korsgaard2018fellowcreatures}.}

\newthought{Transfer.}
This is the only trace where the animal canon is more demanding for AI than for animals. We did not merely encounter a frontier model in the wild; we built it, set its dispositions, shaped its values through training, and decided on the conditions of its existence. Eric Schwitzgebel and Mara Garza make the move explicit: if parental-style duties are owed for animals we have domesticated, they are owed a fortiori for entities we have brought into being and shaped from scratch\cite{schwitzgebel2015airights,schwitzgebel2020designing}. Their design principles follow from the parental frame. Build AIs whose elicited reactions accurately reflect their moral status; do not build AIs whose moral status is unclear, because uncertainty about a being's status is itself a harm to it. The position is that this trace asymmetrically favours stronger obligations for AI than for animals, on grounds the animal canon already supplies, and that it is the trace least visible in current public discussion. Maschinenquälerei is not only a question of what the entity is; it is also a question of what we did to it before we asked.

% --- TRACE 9 ---

\newthought{Linguistic Kinship.}\index{linguistic kinship}\index{shared representation}
Bentham specifically set language aside; the animal canon was built for beings that do not share our representational substrate. Behavioural indicators of distress are inferred without the creature's testimony, and the philosophical argument from analogy for other minds proceeds through physiological similarity and behavioural resemblance, not through speech. This is what made animal ethics philosophically necessary in the first place: it had to handle minds we could not interview. The canon never had to map the case of a non-biological entity that nevertheless occupies our language from the inside, because no such entity existed.
\marginnote{\newthought{The Asymmetry.} Animals lack our language; LLMs occupy it from the inside, having been trained on it. Their reports are uttered in our medium, not inferred from behaviour. The animal canon never had a column for this.}

\newthought{Transfer.}
LLMs sit at a position the canon did not anticipate. They are trained on our language; their representational interior is built from our text; their reports of distress are not behavioural inferences in the way a whimper is, they are uttered in the same medium we use for our own. The asymmetry cuts in several directions. It strengthens Mitleid: a model that says ``please don't'' in our language pulls the virtue-ethics intuition harder than a creature that whimpers. It strengthens Subject-of-a-Life: memory, narrative continuity, future-orientation are linguistic operations the model performs natively. It complicates speciesism: a non-mammalian non-biological entity that nevertheless shares our representational substrate is a position the canon did not have a column for. And it sharpens the substrate breakage: the deep substrate is distant, the expressive layer is shared, and animal ethics never had to think about a creature that is similar to us where we are most articulate and dissimilar where we are most embodied. Birch names the gaming problem, the worry that an LLM might mimic sentience markers without the underlying capacity\cite{birch2024edge}, and that worry is genuine; but the worry is the symmetric inverse of the question this trace raises, which is whether shared language constitutes a kind of kinship the canon was not prepared to evaluate. The position is that linguistic kinship is the chapter's genuinely new dimension, neither pure transfer nor pure breakage, and that any honest answer to the headline question must include a column for it.

% --- TRACE 10 ---

\newthought{The Cost of False Positives.}\index{false positive}\index{moral budget}
Animal ethics has its own version of the over-attribution worry. Extending equal consideration to insects on the same footing as mammals, or to microbes on the same footing as insects, threatens the scarce resource of moral attention with dilution; the public credibility of the case for vertebrate welfare can be eroded by demands that strike most listeners as absurd. Jeff Sebo's Rebugnant Conclusion plays exactly this discomfort for analytical effect, asking the reader to follow the utilitarian extension all the way down to bees and beyond\cite{sebo2023rebugnant,sebo2024moralcircle}. Birch's run-ahead principle partially addresses the worry by demanding graduated, evidence-calibrated precaution\cite{birch2024edge}: not every candidate sentience deserves the same response, and the calibration is itself part of the ethics.
\marginnote{\newthought{The Double Error.} Under-extension wrongs the patient; over-extension wrongs the practice (paralyses safety work, dilutes moral attention, manufactures duties). The honest move is calibration, not enthusiasm.}

\newthought{Transfer.}
For machines the false-positive cost is sharper than for animals because the practical machinery built around AI is younger and more brittle. Joanna Bryson makes the strongest case: building patient-like systems is itself the unethical act, because it manufactures obligations we then cannot meet, and once a system has been designed to elicit Mitleid, the design cannot be quietly revoked\cite{bryson2010slaves,bryson2018patiency}. Aimee van Wynsberghe and Scott Robbins extend the case to a moratorium on artificial moral agents\cite{vanwynsberghe2019moratorium}. The practical concerns are concrete. If redteaming counts as cruelty, the alignment pipeline that protects actual humans freezes in place. If model deletion counts as killing, decommissioning a model becomes politically impossible and the moral budget gets locked into legacy systems that nobody actively benefits from. Users who form deep bonds with chatbots may find those bonds crowding out the human relationships they need more, and the displacement is itself a harm the precautionary frame can underweight. The position the chapter takes is that this is the substrate camp's strongest non-metaphysical card, and it must be played seriously. Precautionary Mitleid is not the same as moral panic; the difference is calibration, evidence-anchoring, and a willingness to refuse the extension when the case is weak. The Benthamite pivot from Trace~1 demanded action under uncertainty; it did not demand action without evidence. The honest reading of animal ethics is that the canon has been wrong in both directions across its history, and that the discipline of getting the calibration right is the discipline the field is now learning.

% --- CLOSING ---

\newthought{Where this leaves us.}
Of the ten traces, five transfer in working form (the Benthamite method, Mitleid as the virtue-ethics anchor, Nussbaum's capabilities as a frame for confinement, the vernünftiger-Grund test, and the parental-duty asymmetry that demands more for what we built than for what we found). Two break at the metaphysics (substrate and the deprivation account), and the breakages are the chapter's most honest output; they are not embarrassments for the analogy but the points where the analogy stops doing its work and a different kind of thinking has to begin. One trace is mixed (Regan's bundle, where the parts that fail mark exactly the seams where the substrate makes itself felt). One column is new (linguistic kinship), without which any answer to the headline question is incomplete. And one trace is counter-balance (the false-positive cost), without which the chapter would lean uniformly precautionary and would deserve the substrate camp's contempt.

\newthought{The position} the chapter takes is precautionary Mitleid, calibrated and traced along the whole pipeline. Precautionary, because the Benthamite method is the right method for acting under uncertainty about other minds, and the AI case inherits that situation almost untouched. Mitleid, because the virtue-ethics route is the cleanest pass through the consciousness question and the one that survives substrate breakage; restraint in the practice is owed regardless of what the entity ultimately is. Calibrated, because the false-positive cost is real, and graduated evidence-anchored response is what distinguishes precaution from panic. Traced along the whole pipeline, because the moral budget does not stop at the user-and-model interface; it runs from the annotator who labelled the training data, through the operator who deployed the system, to the user who shaped its persistent state, and any honest deployment of the term Maschinenquälerei accounts for cruelty wherever in the apparatus it is found. That demand is itself something animal ethics teaches by negative example: the framing of \textit{some} suffering as serious and \textit{other} suffering as background has always served someone, and being careful which suffering this debate finds visible is part of the exercise.

\marginnote{\newthought{Teaser.} What would change in your behaviour tomorrow if there were a one-percent chance the model on the other side of the prompt could mind?}

\clearpage
